\clearpage
\chapter*{Fe de erratas y complemento a la exposición n.º 236}
\addcontentsline{toc}{chapter}{Fe de erratas y complemento a la exposición n.º 236}
\src{303}

\noindent\textbf{8. TDTE VI: Los esquemas de Picard. Propiedades generales.}

\begin{center}
[Seminario Bourbaki, t. 14, 1961/62, n.º 236, 23 pp.]
\end{center}

\noindent\textbf{Página 12, línea 16.} --- En lugar de «recta afín», léase
«recta afín privada del origen».

\noindent\textbf{Línea 18.} --- En lugar de «$X[t]$», léase
«$X[t,t^{-1}]$».

\noindent\textbf{Página 16, proposición 3.1, 6.ª línea.} --- En lugar de «Ahora
la hipótesis ...», léase «Ahora bien, de la hipótesis ...».

\noindent\textbf{Página 21, § 4: El teorema de finitud para el esquema de
Picard.} --- Las cuestiones de finitud del género de las planteadas en este
parágrafo han quedado resueltas casi por completo desde que se redactó la
presente exposición. Indiquemos los principales hechos conocidos actualmente
en esta dirección. Para simplificar los enunciados, suponemos implícitamente
que existen todos los preesquemas de Picard que intervienen en ellos, aunque
una modificación evidente de estos enunciados permite prescindir de toda
hipótesis de existencia. En lo sucesivo, $S$ designa un esquema noetheriano y
$X,Y$ esquemas propios sobre $S$.

\medskip
\noindent(i) Sea $f:X\longrightarrow Y$ un morfismo sobreyectivo. Entonces
\[
f^*: \operatorname{Pic}_{Y/S}\longrightarrow\operatorname{Pic}_{X/S}
\]
es un morfismo de tipo finito.

\src{304}
\noindent(ii) Supongamos que $Y$ es proyectivo sobre $S$ y está provisto de un
Módulo invertible relativamente amplio sobre $S$, y sea $X$ el esquema de ceros
de una sección cualquiera de este Módulo, con $f:X\longrightarrow Y$ la inmersión
canónica. Supongamos, por último, que las componentes irreducibles de las
fibras de $Y/S$ tienen dimensión $\geq3$. Entonces
\[
f^*: \operatorname{Pic}_{Y/S}\longrightarrow\operatorname{Pic}_{X/S}
\]
es de tipo finito.

\noindent(iii) Supongamos que $X$ es proyectivo sobre $S$ y que todas sus
fibras geométricas son íntegras y de dimensión $n$. Sea $\mathcal O_X(1)$ un
Módulo invertible amplio sobre $X$ que permita definir polinomios de Hilbert. Sea
$M$ una parte de $\operatorname{Pic}_{X/S}$; entonces $M$ es cuasicompacta si
y solo si, en los polinomios de Hilbert
\[
a_0x^n+a_1x^{n-1}+a_2x^{n-2}+\cdots
\]
de los elementos de $M$, los coeficientes $a_1$ y $a_2$ permanecen acotados, omitiendo todo coeficiente cuyo monomio tendría grado negativo; así, solo interviene $a_1$ si $n=1$, y ninguno de los dos si $n=0$.

\noindent(iv) Para todo entero $n\neq0$, el homomorfismo de potencia $n$-ésima
del preesquema en grupos $\operatorname{Pic}_{X/S}$ es un morfismo de tipo
finito.

Obsérvese que (i) y (ii) significan también que, bajo las hipótesis
consideradas respectivamente en esos enunciados, una parte $M$ de
$\operatorname{Pic}_{Y/S}$ es cuasicompacta si y solo si lo es su imagen en
$\operatorname{Pic}_{X/S}$. Se deduce que un Módulo invertible $\mathcal L$ sobre
$Y$ es $\tau$-equivalente a cero si y solo si lo es su imagen inversa sobre
$X$; en otros términos, $\operatorname{Pic}_{Y/S}^\tau$ es la imagen inversa
de $\operatorname{Pic}_{X/S}^\tau$. En particular, para demostrar que el
primer preesquema es de tipo finito sobre $S$, basta demostrarlo para el
segundo, puesto que $\operatorname{Pic}_{Y/S}\longrightarrow
\operatorname{Pic}_{X/S}$ es de tipo finito. Utilizando entonces (i), el lema
de Chow y (iii), se obtiene:

\medskip
\noindent(v) $\operatorname{Pic}_{X/S}^\tau$ es de tipo finito sobre $S$.

En general, la conjunción de (i) para un morfismo finito y de (ii) permite
reducir la mayoría de las cuestiones de finitud al caso en que $X/S$ tiene
fibras geométricas íntegras normales de dimensión $\leq2$; a menudo, aplicando
(i) a un morfismo sobreyectivo no necesariamente finito y la resolución de
singularidades de las superficies algebraicas (debida, en característica
arbitraria, a Abhyankar), se reduce incluso al caso en que $X/S$ es liso y,
por tanto, tiene fibras geométricas no singulares de dimensión 2. Esto permite,
por ejemplo, teniendo en cuenta (v) y la desigualdad de Picard--Igusa que
acota el rango del grupo de Néron--Severi de una superficie proyectiva no
singular, demostrar la siguiente generalización del teorema de finitud de
Néron:

\medskip
\noindent(vi) Sea $X/S$ propio sobre $S$ y por lo demás arbitrario. Entonces
los grupos de Néron--Severi
\[
\operatorname{Pic}_{X_i/k_i}/\operatorname{Pic}_{X_i/k_i}^\circ
\]
de las fibras geométricas $X_i/k_i$ de $X/S$

\src{305}
son de tipo finito, y sus rangos y los órdenes de sus subgrupos de torsión
permanecen acotados.

El mismo método de reducción al caso de superficies no singulares, junto con
los teoremas conocidos para este caso (a saber, el teorema de finitud de Néron
y el hecho de que la forma de intersección sobre el grupo de Néron--Severi es
no degenerada), implica:

\medskip
\noindent(vii) Sea $X$ un esquema propio sobre un cuerpo algebraicamente
cerrado. Existe un número finito de curvas cerradas íntegras $C_i$
($1\leq i\leq r$) en $X$ con la propiedad siguiente: para toda parte $M$ de
$\operatorname{Pic}_{X/k}$, $M$ es cuasicompacta si y solo si los enteros
\[
\deg_{C_i'}\mathcal L\qquad(\mathcal L\in M)
\]
permanecen acotados (donde $C_i'$ designa la normalización de $C_i$).

Aquí se puede tomar como $r$ el rango del grupo de Néron--Severi. Una vez
sabido que este último es de tipo finito, (vii) se reduce al hecho de que las
formas lineales sobre el grupo de Néron--Severi definidas por curvas $C$ en
$X$ solo se anulan simultáneamente sobre los elementos de torsión del grupo de
Néron--Severi. En el caso en que $X$ es proyectivo y no singular, este
resultado, así como (v), se debía a Matsusaka. Utilizando (vii), se obtiene
fácilmente la siguiente caracterización de los Módulos invertibles
$\tau$-equivalentes a cero sobre $X$:

\medskip
\noindent(viii) Sean $X/k$ un esquema propio sobre un cuerpo y $\mathcal L$ un
Módulo invertible sobre $X$. Las condiciones siguientes son equivalentes:

\noindent a. $\mathcal L$ es $\tau$-equivalente a cero;

\noindent b. Para todo Módulo coherente $F$, se tiene
\[
\chi(F\otimes\mathcal L)=\chi(F),
\]
donde $\chi$ designa la característica de Euler--Poincaré;

\noindent b'. Como (b), pero con $F=\mathcal O_Y$, donde $Y$ es un
subesquema cerrado íntegro de dimensión 1 en $X$;

\noindent c. Para todo $Y$ como el anterior, si $Y'$ designa la curva
normalizada, se tiene
\[
\deg\mathcal L_{Y'}=0.
\]

Si $X/k$ es proyectivo y está provisto de un Módulo invertible amplio
$\mathcal O_X(1)$, las condiciones precedentes son también equivalentes a las
siguientes:

\noindent d. Para todo entero $m$, $\mathcal L^{\otimes m}(1)$ es amplio.

\noindent e. (Si $X$ es íntegro.) Para todo par de enteros $m,n$, se tiene
\[
\chi(\mathcal L^{\otimes m}(n))=\chi(\mathcal O_X(n))
\]

\src{306}
(es decir, (b) es cierto tomando $F=\mathcal L^{\otimes m}(n)$).

Para la suficiencia de esta última condición, obsérvese que significa que los
polinomios de Hilbert de los $\mathcal L^{\otimes m}$ son todos iguales; por
consiguiente, en virtud del criterio de Mumford (iii), los
$\mathcal L^{\otimes m}$ permanecen en una parte cuasicompacta de
$\operatorname{Pic}_{X/k}$, es decir, se tiene (a). Las condiciones (b), (b')
y (c) deben considerarse variantes, sobre un esquema propio arbitrario, de la
noción de equivalencia numérica definida habitualmente sobre variedades
proyectivas no singulares. Para estas últimas, la equivalencia de (a) y (c)
era, naturalmente, bien conocida (Matsusaka).

El criterio (e) implica asimismo el resultado siguiente:

\medskip
\noindent(ix) Sea $f:X\longrightarrow S$ un morfismo proyectivo y plano cuyas
fibras geométricas son íntegras. Entonces $\operatorname{Pic}_{X/S}^\tau$ es
abierto y cerrado en $\operatorname{Pic}_{X/S}$.

Nos limitaremos a algunos comentarios sobre la demostración de los resultados
clave (i), (ii) y (iii) (el resultado (iv) se aparta un poco de los demás y se
demuestra utilizando únicamente (i) para los morfismos finitos sobreyectivos
radiciales, concretamente para un morfismo de Frobenius). Para (i), se
utilizan de manera esencial las ideas del descenso no plano (véase en
particular TDTE I, página 9). Resulta que, dado que solo se buscan resultados
de finitud, la falta de criterios de efectividad para datos de descenso es
inofensiva. Por otra parte, Mumford ha demostrado recientemente una forma un
poco más débil de (iii), a saber, el criterio que hace intervenir todos los
coeficientes del polinomio de Hilbert. Su argumento es extremadamente simple
y está inspirado en la demostración de un criterio de amplitud de Nakai
(enunciado por este último para las superficies no singulares y generalizado
por Mumford a morfismos proyectivos arbitrarios). Me parece, sin embargo, que
este argumento solo es válido bajo una ligera restricción adicional sobre las
fibras de $X/S$, la propiedad $(S_2)$ de Serre, que se verifica cuando las
fibras geométricas son normales. Se utiliza entonces este criterio restringido
en la demostración de (ii): el criterio (i) permite, en efecto, reducirse al
caso en que $Y/S$ es plano con fibras geométricas íntegras y normales y,
aplicando el criterio de Mumford, se reduce fácilmente al caso en que $X/S$
satisface las mismas condiciones. La hipótesis de dimensión implica entonces
que las fibras geométricas de $Y$ y de $X$ tienen profundidad $\geq2$ en sus
puntos cerrados, lo que permite aplicar los «criterios de equivalencia» en la
forma que se les da en SGA 1962 (exposición XII) y concluir la demostración de
(ii). Una vez que se dispone de (i) y (ii), no es difícil eliminar del criterio
de Mumford toda hipótesis de normalidad sobre las fibras y demostrarlo también
en la forma más fuerte dada en (iii).

\src{307}
Obsérvese, por último, que la demostración de (i) y (ii) muestra asimismo que,
cuando $S$ es el espectro de un cuerpo $k$, el morfismo
\[
\operatorname{Pic}_{Y/k}\longrightarrow\operatorname{Pic}_{X/k}
\]
es afín (y no solamente de tipo finito).
